\documentclass[12pt,a4paper]{article}

% --- فراخوانی بسته استایل قالب طرح پیشنهادی و نامه اداری ---
\usepackage{proposal-style}

% =========================================================================
% متغیرها و تنظیمات مشخصات طرح پیشنهادی و نامه اداری (سفارشی‌سازی)
% =========================================================================
\newcommand{\UniversityName}{دانشگاه صنعتی امیرکبیر}
\newcommand{\UniversitySubtitle}{(پلی‌تکنیک تهران)}
\newcommand{\FacultyName}{دانشکده مهندسی برق}

% مشخصات مقام مخاطب نامه (گیرنده)
\newcommand{\RecipientGender}{جناب آقای} % یا: سرکار خانم
\newcommand{\RecipientName}{دکتر رضا صراف شیرازی}
\newcommand{\RecipientTitle}{ریاست محترم دانشکده مهندسی برق}

% مشخصات تشکل / انجمن / گروه ارائه‌دهنده
\newcommand{\SenderTitle}{انجمن ورزشی دانشجویان دانشکده مهندسی برق}
\newcommand{\SenderCouncilTitle}{شورای انجمن ورزشی دانشکده مهندسی برق}
\newcommand{\SenderShortTitle}{انجمن ورزشی}

% موضوع نامه و عنوان طرح پیشنهادی
\newcommand{\LetterSubject}{درخواست تخصیص فضای اداری اختصاصی به انجمن ورزشی دانشجویان دانشکده مهندسی برق}
\newcommand{\ProposalTitle}{طرح پیشنهادی تخصیص فضای اداری به انجمن ورزشی}
\newcommand{\ProposalDate}{\today}
\newcommand{\ProposalAttachment}{دارد (طرح پیشنهادی)}

% =========================================================================
% تنظیمات سربرگ و پانویس (fancyhdr) برای صفحات دوم به بعد
% =========================================================================
\pagestyle{fancy}
\fancyhf{}
\fancyhead[R]{\small\bfseries \ProposalTitle}
\fancyhead[L]{\small\bfseries \FacultyName}
\renewcommand{\headrulewidth}{0.4pt}
\fancyfoot[C]{\small صفحه \thepage\ از \pageref{LastPage}}

% استایل اختصاصی صفحه اول (کاملاً سفید و بدون هدر و بدون شماره صفحه جهت پاراف و هامش)
\fancypagestyle{firstpage}{
    \fancyhf{}
    \renewcommand{\headrulewidth}{0pt}
    \renewcommand{\footrulewidth}{0pt}
}

\begin{document}
\thispagestyle{firstpage}

% =========================================================================
% سربرگ رسمی نامه اداری (دانشگاه صنعتی امیرکبیر - دانشکده مهندسی برق)
% =========================================================================
\noindent
\begin{minipage}[c]{0.51\textwidth}
    \begin{flushright}
        \begin{tabular}{@{}c@{\hspace{0.15cm}}r@{}}
            \includegraphics[height=1.9cm,keepaspectratio]{aut-logo.png} &
            \begin{tabular}{@{}r@{}}
                {\large\titr \UniversityName} \\[0.06cm]
                {\small\bfseries \UniversitySubtitle} \\[0.06cm]
                {\normalsize\bfseries \FacultyName} \\[0.06cm]
                {\small\bfseries \SenderTitle}
            \end{tabular}
        \end{tabular}
    \end{flushright}
\end{minipage}%
\hfill
\begin{minipage}[c]{0.11\textwidth}
    \begin{center}
        \vspace{-0.3cm}
        {\large\bfseries بسمه تعالی}
    \end{center}
\end{minipage}%
\hfill
\begin{minipage}[c]{0.36\textwidth}
    \begin{flushleft}
        \begin{tabular}{@{}l@{\hspace{0.15cm}}c@{}}
            \begin{tabular}{@{}r@{}}
                {\small\bfseries تاریخ: \lr{\ProposalDate}} \\[0.15cm]
                {\small\bfseries پیوست: \ProposalAttachment}
            \end{tabular} &
            \includegraphics[height=1.7cm,keepaspectratio]{EE-aut.jpg}
        \end{tabular}
    \end{flushleft}
\end{minipage}

\vspace{0.25cm}
\noindent\rule{\textwidth}{1pt}
\vspace{0.6cm}

% =========================================================================
% نامه اداری روکش (صفحه اول - جهت بررسی، دستور و پاراف ریاست محترم دانشکده)
% =========================================================================
{\fontsize{13pt}{24pt}\selectfont
\noindent
\textbf{\RecipientGender{} \RecipientName} \par
\textbf{\RecipientTitle} \par\vspace{0.45cm}
\textbf{موضوع:} \LetterSubject \par\vspace{0.45cm}
\textbf{با سلام و احترام؛} \par\vspace{0.35cm}
بدین‌وسیله به استحضار می‌رساند، \SenderTitle{} با هدف ارتقای نشاط، سلامت روحی و جسمی دانشجویان و اعتلای جایگاه ورزشی دانشکده در سطح مسابقات و المپیادهای دانشگاه صنعتی امیرکبیر، اقدام به تدوین \textbf{«طرح پیشنهادی تخصیص فضا و ساماندهی ورزش دانشجویی»} نموده است.

\vspace{0.35cm}
تحقق اهداف این طرح و خروج فعالیت‌های ورزشی از حالت پراکنده و مقطعی، نیازمند برخورداری انجمن از حداقل زیرساخت‌های اجرایی و تخصیص یک فضای فیزیکی مستقل و مستمر (اتاق اختصاصی) در سطح دانشکده است؛ این فضا به عنوان پایگاهی ثابت برای برگزاری جلسات هماهنگی شورا و تیم‌ها، آنالیزهای فنی، ساماندهی لیگ‌های درون‌دانشکده‌ای، پاسخ‌گویی به مراجعین و نگهداری تجهیزات سبک مورد استفاده قرار خواهد گرفت.

\vspace{0.35cm}
نظر به آغاز رویدادها و المپیادهای دانشگاهی، \textbf{متن کامل طرح پیشنهادی (شامل ضرورت‌ها، برنامه‌های عملیاتی و خدمات مورد نظر) به پیوست این نامه تقدیم حضور می‌گردد.}

\vspace{0.35cm}
خواهشمند است با توجه به رویکرد حمایتی همیشگی جنابعالی از نشاط و فعالیت‌های دانشجویی، با تخصیص یک اتاق به این مجموعه موافقت فرموده و دستورات مقتضی را جهت هماهنگی و اقدام ابلاغ فرمایید.

\vspace{0.35cm}
پیشاپیش از بذل توجه، حسن نظر و اوامر مساعد جنابعالی صمیمانه کمال تشکر و امتنان را دارد.

\vspace{0.9cm}
\begin{flushleft}
    \textbf{با احترام و سپاس،} \par\vspace{0.2cm}
    \textbf{\SenderCouncilTitle}
\end{flushleft}
}

% =========================================================================
% صفحه دوم به بعد: متن کامل طرح پیشنهادی (پیوست نامه اداری)
% =========================================================================
\newpage

\begin{center}
    {\large\bfseries طرح پیشنهادی پیوست} \\[0.25cm]
    {\LARGE\titr \ProposalTitle{} دانشکده \par}
    \vspace{0.2cm}
    {\large\bfseries \FacultyName{} ، \UniversityName \par}
    \vspace{0.3cm}
\end{center}

\noindent
\textbf{ارائه‌دهنده:} \SenderTitle \\
\textbf{مخاطب:} ریاست محترم و اعضای محترم شورای \FacultyName \\
\textbf{موضوع:} \LetterSubject \\

% =========================================================================
% ۱. مقدمه و بیان مسئله
% =========================================================================
\section{مقدمه و بیان مسئله}
ورزش دانشجویی یکی از مؤثرترین بسترها برای افزایش نشاط و پویایی محیط دانشگاه، تقویت روحیه همکاری و کار گروهی، ایجاد تعامل میان دانشجویان و شناسایی و پرورش استعدادهای فردی و جمعی است. حضور فعال و سازمان‌یافته دانشجویان در فعالیت‌های ورزشی، علاوه بر آثار مثبت فردی، می‌تواند به تقویت هویت جمعی و افزایش اعتبار و جایگاه دانشکده در سطح دانشگاه منجر شود.

در این راستا، \SenderTitle{} با هدف گسترش فرهنگ ورزش، افزایش مشارکت دانشجویان در فعالیت‌های ورزشی، شناسایی استعدادها، تشکیل و پشتیبانی از تیم‌های ورزشی و ایجاد بستری مناسب برای حضور قدرتمند دانشکده در رقابت‌های دانشگاهی، درصدد اجرای طرحی جامع و بلندمدت برای توسعه ورزش دانشجویی است.

در حال حاضر، نبود یک فضای مشخص برای استقرار انجمن، برگزاری جلسات، هماهنگی میان تیم‌های ورزشی، پیگیری امور اجرایی و ارتباط مستقیم با دانشجویان، روند سازماندهی و اجرای برنامه‌ها را با محدودیت مواجه می‌کند. پراکندگی این فعالیت‌ها در فضاهای مختلف، امکان پیگیری منظم امور و ایجاد یک ساختار اجرایی پایدار را کاهش می‌دهد. ازاین‌رو، اختصاص یک فضای اداری مشخص می‌تواند بستر لازم برای ساماندهی فعالیت‌ها و افزایش اثربخشی عملکرد انجمن را فراهم سازد.

% =========================================================================
% ۲. اهداف اصلی طرح
% =========================================================================
\section{اهداف اصلی طرح}
این طرح با اهداف زیر تدوین شده است:
\begin{enumerate}
    \item ایجاد ساختاری منسجم، منظم و پایدار برای مدیریت و توسعه فعالیت‌های ورزشی دانشکده.
    \item افزایش مشارکت دانشجویان در فعالیت‌های ورزشی و فراهم کردن فرصت حضور دانشجویان علاقه‌مند در رشته‌های مختلف.
    \item شناسایی استعدادهای ورزشی و تشکیل تیم‌های توانمند برای نمایندگی دانشکده در رقابت‌های دانشگاهی.
    \item ارتقای سطح فنی، آمادگی و هماهنگی تیم‌های ورزشی از طریق برنامه‌ریزی، جلسات تخصصی و بررسی عملکرد.
    \item توسعه مسابقات و رویدادهای ورزشی درون‌دانشکده‌ای و ایجاد فضای رقابتی سالم میان دانشجویان.
    \item حضور شایسته و مؤثر تیم‌های دانشکده در مسابقات و المپیادهای ورزشی دانشگاهی.
    \item ایجاد بستری ماندگار برای استمرار و توسعه فعالیت‌های انجمن در سال‌های تحصیلی آینده.
\end{enumerate}

% =========================================================================
% ۳. محورهای اجرایی و برنامه‌های پیشنهادی
% =========================================================================
\section{محورهای اجرایی و برنامه‌های پیشنهادی}

\subsection*{الف) سازماندهی و تشکیل تیم‌های ورزشی دانشکده}
یکی از مهم‌ترین وظایف انجمن، شناسایی دانشجویان علاقه‌مند و مستعد در رشته‌های مختلف ورزشی و ایجاد ساختاری منظم برای عضوگیری و تشکیل تیم‌های دانشکده است. اقدامات پیش‌بینی‌شده در این بخش عبارت‌اند از:
\begin{itemize}
    \item شناسایی، جذب و استعدادیابی دانشجویان علاقه‌مند در تمامی مقاطع تحصیلی.
    \item تشکیل و ساماندهی تیم‌های ورزشی دانشکده در رشته‌های متنوع انفرادی و تیمی.
    \item تعیین سرپرستان رشته‌ها و ثبت مشخصات ورزشکاران دانشکده.
    \item ایجاد بستر ارتباطی منظم و پایدار میان اعضای تیم‌ها و انجمن ورزشی.
\end{itemize}

\subsection*{ب) پشتیبانی، هماهنگی فنی و برنامه‌ریزی برای مسابقات دانشگاهی}
با توجه به تقویم مسابقات و المپیادهای ورزشی دانشگاهی، ایجاد آمادگی و پیگیری امور اجرایی تیم‌ها نقش تعیین‌کننده‌ای در عملکرد آنان دارد. در همین راستا، اقدامات زیر در دستور کار قرار خواهد گرفت:
\begin{itemize}
    \item برگزاری جلسات هماهنگی اعضای تیم‌ها پیش از آغاز مسابقات.
    \item پیگیری امور اداری، ثبت‌نام و هماهنگی‌های لازم برای حضور تیم‌ها در رقابت‌ها.
    \item بررسی ترکیب تیم‌ها و تبادل نظر پیرامون مسائل اجرایی و فنی.
    \item بررسی عملکرد تیم‌های دانشکده و تحلیل نقاط قوت و زمینه‌های قابل بهبود.
    \item پیگیری تدارکات و ملزومات اولیه موردنیاز برای اعزام و پشتیبانی از تیم‌ها.
\end{itemize}
هدف از این اقدامات، فراهم کردن شرایطی است که تیم‌های دانشکده بتوانند با آمادگی، انسجام و روحیه تیمی بالاتری در رقابت‌ها حضور یافته و نماینده‌ای شایسته برای دانشکده مهندسی برق باشند.

\subsection*{ج) طراحی و برگزاری مسابقات درون‌دانشکده‌ای}
برگزاری مسابقات داخلی می‌تواند نقش مهمی در افزایش مشارکت دانشجویان، شناسایی استعدادهای جدید و ایجاد شور و نشاط در محیط دانشکده داشته باشد. در این بخش، انجمن ورزشی برنامه‌ریزی و اجرای رویدادهای ورزشی مختلف را در دستور کار قرار خواهد داد؛ از جمله:
\begin{itemize}
    \item طراحی و برگزاری مسابقات و لیگ‌های درون‌دانشکده‌ای در رشته‌های مختلف.
    \item هماهنگی با مجموعه تربیت‌بدنی دانشگاه برای دریافت و رزرو سانس سالن‌های ورزشی و زمین چمن.
    \item تنظیم برنامه زمانی مسابقات و هماهنگی با تیم‌های شرکت‌کننده.
    \item اطلاع‌رسانی و ثبت‌نام دانشجویان علاقه‌مند.
    \item پیگیری امور اجرایی و نظارت بر روند برگزاری مسابقات.
    \item ایجاد فرصت رقابت سالم و تعامل بیشتر میان دانشجویان رشته‌ها و ورودی‌های مختلف.
\end{itemize}

\subsection*{د) توسعه فعالیت‌های ورزشی و مشارکت دانشجویان}
انجمن ورزشی علاوه بر فعالیت‌های رقابتی، به دنبال ایجاد بستری برای گسترش فعالیت‌های عمومی ورزشی در دانشکده است. از جمله برنامه‌های این بخش می‌توان به موارد زیر اشاره کرد:
\begin{itemize}
    \item شناسایی علایق و نیازهای ورزشی دانشجویان.
    \item معرفی رشته‌های ورزشی مختلف و فراهم کردن زمینه مشارکت دانشجویان.
    \item برگزاری برنامه‌های ورزشی جمعی و رویدادهای مناسبتی.
    \item تشویق دانشجویان به فعالیت‌های ورزشی و تعاملات گروهی.
    \item ایجاد ارتباط مؤثر میان دانشجویان، ورزشکاران و انجمن ورزشی.
\end{itemize}

\subsection*{هـ) مستندسازی و ایجاد ساختار پایدار}
به‌منظور جلوگیری از مقطعی شدن فعالیت‌ها و حفظ تجربیات به‌دست‌آمده، انجمن در نظر دارد فرآیند مستندسازی و سازماندهی امور را نیز دنبال کند. این اقدامات شامل تشکیل بانک اطلاعاتی ورزشکاران، ثبت نتایج مسابقات، نگهداری سوابق تیم‌ها، تدوین تقویم سالانه فعالیت‌های ورزشی و ثبت تجربیات اجرایی برای استفاده در دوره‌های بعدی انجمن خواهد بود.

% =========================================================================
% ۴. ضرورت تخصیص فضای اختصاصی به انجمن ورزشی
% =========================================================================
\section{ضرورت تخصیص فضای اختصاصی به انجمن ورزشی}
اجرای منظم و مؤثر برنامه‌های فوق، نیازمند برخورداری از یک فضای مشخص و در دسترس است که به‌عنوان پایگاه اجرایی و هماهنگی انجمن مورد استفاده قرار گیرد. این فضا صرفاً محل استقرار اعضای انجمن نخواهد بود، بلکه کارکردهای مختلفی در فرآیند توسعه ورزش دانشکده خواهد داشت، از جمله:
\begin{itemize}
    \item برگزاری جلسات رسمی و برنامه‌ریزی فعالیت‌های انجمن.
    \item برگزاری جلسات فنی و هماهنگی اعضای تیم‌های ورزشی.
    \item انجام بررسی‌ها و آنالیزهای فنی مسابقات.
    \item برنامه‌ریزی برای ترکیب تیم‌ها و هماهنگی بازیکنان.
    \item پیگیری امور مربوط به رزرو سالن‌های ورزشی و زمین چمن.
    \item تنظیم برنامه مسابقات درون‌دانشکده‌ای و هماهنگی با تیم‌های شرکت‌کننده.
    \item ثبت‌نام و سازماندهی دانشجویان علاقه‌مند به فعالیت‌های ورزشی.
    \item نگهداری اسناد، سوابق، اطلاعات و تجهیزات سبک موردنیاز انجمن.
\end{itemize}
همچنین وجود یک فضای ثابت، امکان مراجعه مستقیم دانشجویان برای طرح پیشنهادها، بیان مسائل و پیگیری امور ورزشی را فراهم می‌کند و موجب تقویت ارتباط میان انجمن و بدنه دانشجویی دانشکده خواهد شد.

نکته حائز اهمیت آن است که نیاز به این فضا محدود به مسابقات و المپیادهای پیش‌رو نیست؛ بلکه فعالیت‌های انجمن ماهیتی مستمر و سالانه دارند و برای تحقق اهداف بلندمدت خود، به یک پایگاه ثابت و قابل اتکا نیازمندند.

\subsection*{ویژگی‌ها و نحوه استفاده از فضای مورد درخواست}
فضای مورد درخواست، یک اتاق اداری با ظرفیت مناسب برای برگزاری جلسات اعضای انجمن و هماهنگی‌های تیم‌های ورزشی است که امکان استفاده مستمر در طول سال تحصیلی را داشته باشد. این فضا علاوه بر فراهم کردن امکان انجام امور اجرایی انجمن، به‌عنوان محل مراجعه دانشجویان برای ثبت‌نام در برنامه‌های ورزشی، ارائه پیشنهادها و پیگیری مسائل مرتبط با فعالیت‌های ورزشی دانشکده مورد استفاده قرار خواهد گرفت. همچنین امکان نگهداری اسناد، سوابق و تجهیزات سبک انجمن در این فضا فراهم خواهد بود.

% =========================================================================
% ۵. نقش انجمن ورزشی در توسعه فعالیت‌های دانشکده
% =========================================================================
\section{نقش انجمن ورزشی در توسعه فعالیت‌های دانشکده}
انجمن ورزشی به عنوان بازوی اجرایی دانشجویی در دانشکده مهندسی برق، با استقرار در این فضا، مسئولیت‌های عملیاتی زیر را جهت ساماندهی و پویایی فعالیت‌ها بر عهده خواهد گرفت:
\begin{enumerate}
    \item ساماندهی و پیگیری مستمر امور تیم‌های ورزشی دانشکده.
    \item ایجاد ارتباط مؤثر میان دانشجویان علاقه‌مند، ورزشکاران و مسئولان دانشکده.
    \item برنامه‌ریزی و اجرای مسابقات و رویدادهای ورزشی درون‌دانشکده‌ای.
    \item پیگیری امور مربوط به حضور شایسته تیم‌ها در مسابقات و المپیادهای دانشگاهی.
    \item ثبت و مستندسازی فعالیت‌ها، نتایج مسابقات و سوابق تیم‌ها.
    \item دریافت پیشنهادها، پاسخ‌گویی به مراجعین و شناسایی نیازهای ورزشی دانشجویان.
\end{enumerate}

% =========================================================================
% ۶. دستاوردهای مورد انتظار
% =========================================================================
\section{دستاوردهای مورد انتظار}
تحقق این طرح و فراهم‌آمدن پایگاه استقرار انجمن، دستاوردهای ملموس زیر را برای جامعه دانشکده مهندسی برق به همراه خواهد داشت:
\begin{enumerate}
    \item ارتقای نشاط عمومی، سلامت جسمی و بهبود روحیه و انگیزه تحصیلی دانشجویان.
    \item افزایش مشارکت عمومی دانشجویان در فعالیت‌های ورزشی و غنی‌سازی اوقات فراغت.
    \item تقویت روحیه همدلی، تعامل سازنده و هویت جمعی در میان دانشجویان گرایش‌ها و مقاطع مختلف.
    \item ارتقای کیفیت حضور و انسجام تیم‌های دانشکده در مسابقات دانشگاه صنعتی امیرکبیر.
    \item ایجاد زیرساختی ماندگار و انباشت تجربیات برای تداوم پایدار فعالیت‌ها در سال‌های تحصیلی آینده.
\end{enumerate}

% =========================================================================
% ۷. حمایت‌ها و امکانات مورد درخواست
% =========================================================================
\section{حمایت‌ها و امکانات مورد درخواست}
با توجه به اهداف و برنامه‌های مطرح‌شده، انجمن ورزشی دانشکده مهندسی برق تقاضا دارد موارد زیر از سوی ریاست محترم و اعضای محترم شورای دانشکده مورد بررسی و مساعدت قرار گیرد:

\subsection*{الف) درخواست اصلی}
\begin{itemize}
    \item \textbf{تخصیص یک اتاق اداری اختصاصی} به انجمن ورزشی دانشکده به‌عنوان محل استقرار، برنامه‌ریزی، برگزاری جلسات و پیگیری امور اجرایی.
\end{itemize}

\subsection*{ب) درخواست‌های تکمیلی}
\begin{itemize}
    \item تسهیل هماهنگی با مجموعه تربیت‌بدنی دانشگاه برای استفاده از امکانات ورزشی (سالن‌ها و زمین چمن).
    \item حمایت اجرایی از برگزاری مسابقات و رویدادهای ورزشی درون‌دانشکده‌ای.
    \item مساعدت در تأمین تجهیزات اولیه و ملزومات اداری موردنیاز.
    \item حمایت از تیم‌های ورزشی دانشکده برای حضور در مسابقات و المپیادهای دانشگاهی.
\end{itemize}

% =========================================================================
% ۸. تعهدات انجمن در قبال فضای اختصاص‌یافته
% =========================================================================
\section{تعهدات انجمن در قبال فضای اختصاص‌یافته}
انجمن ورزشی دانشکده مهندسی برق در راستای استفاده صحیح و مؤثر از فضای اختصاص‌یافته، خود را متعهد به رعایت موارد زیر می‌داند:
\begin{enumerate}
    \item استفاده از فضای تخصیص‌یافته صرفاً در راستای اهداف و فعالیت‌های انجمن ورزشی دانشکده.
    \item حفظ، نگهداری و مراقبت از فضای فیزیکی و تجهیزات موجود و رعایت مقررات مربوط به استفاده از آن.
    \item تعیین مسئول مشخص از میان اعضای انجمن برای مدیریت امور مربوط به اتاق و نظارت بر نحوه استفاده از آن.
    \item فراهم کردن امکان مراجعه دانشجویان و پاسخ‌گویی به درخواست‌ها و مسائل مرتبط با فعالیت‌های ورزشی دانشکده در زمان‌های مشخص.
    \item ارائه گزارش دوره‌ای از فعالیت‌ها، برنامه‌ها و عملکرد انجمن به ریاست دانشکده یا مسئول مربوطه.
    \item مستندسازی فعالیت‌ها و انتقال اطلاعات و تجربیات به اعضای دوره‌های بعدی انجمن، به‌منظور حفظ استمرار فعالیت‌ها.
\end{enumerate}

% =========================================================================
% ۹. جمع‌بندی و درخواست نهایی
% =========================================================================
\section{جمع‌بندی و درخواست نهایی}
با توجه به موارد مطرح‌شده، تخصیص یک فضای اداری مشخص و مستمر به انجمن ورزشی دانشکده، گامی مؤثر در جهت ساماندهی فعالیت‌های ورزشی، تقویت ارتباط با بدنه دانشجویی و ایجاد بستری پایدار برای اجرای برنامه‌های ورزشی خواهد بود.

انجمن ورزشی دانشکده نیز متعهد است از فضای اختصاص‌یافته در راستای اهداف تعیین‌شده استفاده نموده و با رعایت ضوابط دانشکده، مسئولیت حفظ و نگهداری آن و ارائه گزارش عملکرد فعالیت‌ها را بر عهده بگیرد.

از ریاست محترم و اعضای محترم شورای دانشکده مهندسی برق درخواست می‌شود با بررسی این طرح و موافقت با تخصیص یک فضای اداری اختصاصی به انجمن ورزشی، زمینه اجرای برنامه‌های پیشنهادی و توسعه فعالیت‌های ورزشی دانشکده را فراهم آورند.

\vspace{0.8cm}
\begin{flushleft}
    \textbf{با احترام،} \\[0.2cm]
    \textbf{\SenderCouncilTitle}
\end{flushleft}

\end{document}
